% !TeX root = ../main.tex
% =================================================
% Chapter 6: Conclusions
% =================================================

\chapter{Conclusions} % LaTeX chapter title
\label{ch:conclusions}

Restate the overarching research question from your introduction — not the 
individual paper questions, but the thesis-level one. This closes the loop for 
the reader.

Stochastic Inertia so far only proven for unstable fixed points with a single
unstable dimension. Therefore questionable whether theory is applicable
to sppt in wcbs. Local dimensionality of the weather system is an active
area of research (cite Faranda stuff here). Connect to turbulence theory
ventilated in the intro.

\section{Synthesis}\label{sec:synth}
This is the core and what distinguishes a thesis conclusion from a paper
conclusion. Don't summarise each paper in turn — that reads like an abstract
collection. Instead, organise around themes or findings that cut across the
papers. Where do they agree, reinforce, or tension each other? What emerges
from the combination that wasn't visible in any single paper?

% Limitations
Be honest but concise. Focus on limitations of the body of work as a whole, not
individual paper limitations which are already discussed in each chapter.

\section{Implications and future directions}\label{sec:Imp_Fut}
Broader implications (~2 pages)
Step back from your specific methods and datasets. What does this mean for
the field? Are there methodological lessons, not just scientific ones? This is
where you can be a little bolder than you would in a paper.
Limitations (~1 page)
Be honest but concise. Focus on limitations of the body of work as a whole,
not individual paper limitations which are already discussed in each chapter.
Future directions (~2 pages)
The most valuable part for other researchers. Distinguish between:

Near-term — obvious extensions of what you did
Long-term — bigger open questions your work points toward but couldn't address

Closing paragraph
One strong paragraph that lands the thesis. What is the single most important
thing you want the reader to take away? Write it plainly and confidently.